% ch4.tex -- Module 5, Chapter 4: Gravity, orbits and Kepler's laws
% Dot points: M5.IQ3.01, M5.IQ3.02, M5.IQ3.03, M5.IQ3.04 (energy in orbits: Chapter 5)
% Figures: TikZ inline, plus figures/ch4_cannonball.py, ch4_equalareas.py,
%          ch4_ellipseforces.py (matplotlib, PDF next to each script)
\chapter{Gravity, Orbits and Kepler's Laws}\label{ch4:chap}
\covers{M5.IQ3.01, M5.IQ3.02, M5.IQ3.03, M5.IQ3.04}

\inquiry{How does the force of gravity determine the motion of planets and satellites?}

\begin{goals}
  \item calculate the gravitational force between two masses, and the field strength $g$ on a
        planet and above it
  \item explain why a satellite stays in orbit, and calculate its speed, period and radius
  \item state Kepler's three laws and explain them with forces and energy
  \item compare low Earth orbits and geostationary orbits, and link them to their uses
\end{goals}

\begin{syllabus}
  \item M5.IQ3.01: use Newton's law of universal gravitation to find the force between two
        masses, investigate what affects $g = GM/r^2$, and predict $g$ at any point
  \item M5.IQ3.02: investigate the orbits of planets and satellites, using gravitational force,
        centripetal force and acceleration, mass, orbital radius, orbital velocity and period
  \item M5.IQ3.03: predict the properties of low Earth orbits and geostationary orbits, and link
        them to their uses
  \item M5.IQ3.04: investigate how Kepler's laws are related to the forces on planets and to
        their total energy (energy: \cref{ch5:chap}), using $v = 2\pi r/T$ and
        $r^3/T^2 = GM/4\pi^2$
\end{syllabus}

%% ==================================================================
\section{Newton's law of gravitation and field strength}\label{ch4:sec:newton}
\covers{M5.IQ3.01}

\begin{bigidea}
Every mass pulls every other mass. The pull is bigger for bigger masses, and it gets weaker
quickly with distance. A planet's field strength $g$ is its pull on each kilogram.
\end{bigidea}

\begin{keywords}
  \kw{Gravitational force}{the pull between any two objects that have mass. Unit: newton (N).}
  \kw{Universal gravitational constant ($G$)}{a fixed number that sets how strong gravity is.}
  \kw{Inverse-square law}{a rule where doubling the distance makes the effect four times smaller.}
  \kw{Gravitational field strength ($g$)}{the gravitational force on each kilogram at a point.}
  \kw{Altitude ($h$)}{the height above a planet's surface.}
\end{keywords}

A pen falls because the Earth pulls it. The Moon orbits the Earth for the same reason. In 1687
Isaac Newton showed that one law explains both. It is called the law of \term{universal
gravitation}: ``universal'' means it works everywhere, on Earth and in space. Gravity acts between \emph{any}
two masses, and it always attracts.

\begin{refresher}{Mass, weight and Newton's third law}
\term{Mass} ($m$, in kg) is the amount of matter in an object. It is the same everywhere.
\term{Weight} is the gravitational force of a planet on an object (in N). Near the Earth's
surface, weight $= mg = m \times 9.8$. Newton's third law: if A pulls B, then B pulls A with a
force of the same size in the opposite direction.
\end{refresher}

\begin{sidebyside}[0.5]
  \centering
  \begin{tikzpicture}
    \fill[object] (0,0) circle (0.85);
    \node[diagramlabel] at (0,0.32) {$M$};
    \fill[object] (4.0,0) circle (0.32);
    \node[diagramlabel] at (4.0,0.6) {$m$};
    \fill[hscink] (0,0) circle (1.3pt);
    \fill[hscink] (4.0,0) circle (1.3pt);
    \draw[force] (0.85,0) -- (1.85,0) node[above, pos=0.5, text=hscforce, font=\small] {$F$};
    \draw[force] (3.68,0) -- (2.68,0) node[above, pos=0.5, text=hscforce, font=\small] {$F$};
    \draw[hscmuted, line width=0.6pt, {Stealth[length=2mm]}-{Stealth[length=2mm]}]
      (0,-1.15) -- node[below, diagramlabel] {$r$ (centre to centre)} (4.0,-1.15);
    \draw[guide] (0,-0.85) -- (0,-1.3);
    \draw[guide] (4.0,-0.32) -- (4.0,-1.3);
  \end{tikzpicture}
  \captionof{figure}{Two masses pull each other with forces of equal size.}\label{ch4:fig:pair}
\nextside
  The force depends on (\cref{ch4:fig:pair}):
  \begin{itemize}
    \item the two masses, $M$ and $m$;
    \item the distance $r$ between their \term{centres of mass} (the points where all their
          mass seems to act; for a planet, its centre).
  \end{itemize}
  The arrows have the same length: $M$ pulls $m$ exactly as hard as $m$ pulls $M$.
\end{sidebyside}

\begin{toolkit}{Powers of ten, inverse squares and significant figures}
\begin{itemize}
  \item Type $6.67\times10^{-11}$ with the \texttt{EXP} (or $\times10^x$) key. Put brackets round a
        number before squaring it: $(7.221\times10^{6})^2 = 5.214\times10^{13}$. The power of ten
        doubles.
  \item If $F \propto \dfrac{1}{r^2}$ and $r$ is multiplied by $k$, then $F$ is divided by $k^2$:
        $r \times 2 \Rightarrow F \div 4$, \ $r \times 3 \Rightarrow F \div 9$.
  \item Significant figures: the data sheet gives $M_\mathrm{E}$ to only 2 significant figures.
        HSC answers with 2 or 3 significant figures are usually accepted. In Earth-orbit examples
        this book gives 3, so you can check your working.
\end{itemize}
\end{toolkit}

\begin{formulacard}{Newton's law of universal gravitation}
  \cardformula{F = \frac{GMm}{r^2}}
  \begin{symbols}
    \sym{$F$}{gravitational force on each of the two objects}{\si{N}}
    \sym{$G$}{universal gravitational constant, \num{6.67e-11} (data sheet)}{\si{N.m^2.kg^{-2}}}
    \sym{$M$, $m$}{masses of the two objects ($M$ is usually the bigger one)}{\si{kg}}
    \sym{$r$}{distance between the centres of the two objects}{\si{m}}
  \end{symbols}
  \inwords{the force is bigger when either mass is bigger. It is smaller when the objects are
    further apart: double the distance and the force is 4 times smaller.}
  \whentouse{two point masses or two spheres, with $r$ from centre to centre (planets, moons,
    satellites). For other objects, such as people or cars, it gives only a rough estimate.}
  \whennottouse{inside a planet ($r$ smaller than its radius). This is not in the HSC course.}
  \rearranged{$r = \sqrt{\dfrac{GMm}{F}}$ \qquad $M = \dfrac{Fr^2}{Gm}$}
  \tensecond{two lead spheres of \qty{2.0e3}{kg} and \qty{3.0e3}{kg}, with centres \qty{2.0}{m}
    apart: $F = \dfrac{\num{6.67e-11} \times \num{2.0e3} \times \num{3.0e3}}{2.0^2} = \qty{1.0e-4}{N}$,
    far too small to feel.}
  \trap{$r$ goes from centre to centre, not surface to surface. Square all of $r$, including the
    power of ten.}
\end{formulacard}


\subsection{Gravitational field strength}
A skydiver feels the Earth's pull without touching the Earth. The Earth makes a
\term{gravitational field}: a region where any mass feels a gravitational force. The
\term{gravitational field strength} is the force on each kilogram, $g = \dfrac{F}{m}$.
Its unit is \si{N.kg^{-1}}. This is the same as \si{m.s^{-2}}, because
\qty{1}{N} = \qty{1}{kg.m.s^{-2}}. Newton's second law says $a = \dfrac{F}{m}$. So for an
object in \term{free fall} (gravity is the only force on it), the force on each kilogram is also
its acceleration.

\begin{sidebyside}[0.56]
  \centering
  \begin{tikzpicture}[scale=0.82]
    \fill[hscbrand!10, draw=hscink, line width=0.8pt] (0,0) circle (1.0);
    \node[diagramlabel] at (0,-0.15) {planet};
    \foreach \a in {0,30,...,330} \draw[field] (\a:2.0) -- (\a:1.05);
    \draw[hscmuted, line width=0.6pt] (-0.22,0.82) rectangle (0.22,1.26);
    \node[diagramlabel, align=center] at (0,-2.5) {radial field};
    \begin{scope}[xshift=4.0cm]
      \draw[hscmuted, line width=0.6pt] (-1.1,-1.1) rectangle (1.1,1.5);
      \fill[ground] (-1.1,-1.1) rectangle (1.1,-0.75);
      \draw[hscink, line width=0.8pt] (-1.1,-0.75) -- (1.1,-0.75);
      \foreach \x in {-0.8,-0.4,0,0.4,0.8} \draw[field] (\x,1.3) -- (\x,-0.68);
      \node[diagramlabel, align=center] at (0,-1.65) {near the surface:\\nearly uniform};
    \end{scope}
    \draw[guide] (0.22,1.26) -- (2.9,1.5);
    \draw[guide] (0.22,0.82) -- (2.9,-1.1);
  \end{tikzpicture}
  \captionof{figure}{A planet's field. Arrows show the force on a mass.}\label{ch4:fig:field}
\nextside
  \term{Field lines} are arrows that show which way a mass would be pulled. Around a planet the
  field lines point to the centre. This is a \term{radial field} (\cref{ch4:fig:field}). The
  lines spread out further away, so the field gets weaker.

  \smallskip
  In a small region near the surface the lines are almost parallel. The field there is nearly a
  \term{uniform field} (same size and direction everywhere). That is why projectile questions use
  one value, $g = \qty{9.8}{m.s^{-2}}$.
\end{sidebyside}

Put Newton's law into $g = \dfrac{F}{m}$:
\[ g = \frac{F}{m} = \frac{GMm}{r^2} \div m = \frac{GM}{r^2} \]
The object's mass $m$ cancels. So \emph{all objects fall with the same acceleration} in the same
place, if air resistance is small enough to ignore. In 1971 an astronaut on the airless Moon
dropped a hammer and a feather, and they landed together.

\begin{sidebyside}[0.42]
  \centering
  \begin{tikzpicture}
    \fill[hscbrand!10, draw=hscink, line width=0.8pt] (0,0) circle (1.15);
    \draw[hscmuted!70, dashed, line width=0.6pt] (40:2.55) arc (40:-26:2.55);
    \fill[hscink] (0,0) circle (1.4pt);
    \node[diagramlabel, left, inner sep=2pt] at (0,0) {centre};
    \draw[hscink, line width=0.7pt] (0,0) -- (2.55,0);
    \draw[hscink, line width=0.7pt] (1.15,0.1) -- (1.15,-0.1);
    \node[diagramlabel, above] at (0.62,0) {$r_\mathrm{E}$};
    \node[diagramlabel, above] at (1.85,0) {$h$};
    \draw[guide] (0,0) -- (0,-1.6);
    \draw[guide] (2.55,0) -- (2.55,-1.6);
    \draw[hscmuted, line width=0.6pt, {Stealth[length=2mm]}-{Stealth[length=2mm]}]
      (0,-1.45) -- node[below, diagramlabel, inner sep=4pt] {$r = r_\mathrm{E} + h$} (2.55,-1.45);
    \fill[object] (2.55,0) circle (4.5pt);
    \node[diagramlabel, right, inner sep=5pt] at (2.55,0) {satellite};
    \node[diagramlabel, text=hscmuted] at (0.1,0.75) {Earth};
  \end{tikzpicture}
  \captionof{figure}{Altitude $h$ and distance $r$ from the centre.}\label{ch4:fig:altitude}
\nextside
  Questions usually give the \term{altitude} $h$: the height above the surface. Gravity formulas
  need $r$, measured from the \emph{centre} (\cref{ch4:fig:altitude}). So add the Earth's radius
  to the altitude:
  \[ r = r_\mathrm{E} + h \]
  The data sheet gives the Earth's radius, $r_\mathrm{E} = \qty{6.371e6}{m}$, and its mass,
  $M_\mathrm{E} = \qty{6.0e24}{kg}$.
\end{sidebyside}

\begin{formulacard}[notonsheet]{Distance from the centre and altitude}
  \cardformula{r = r_\mathrm{E} + h}
  \begin{symbols}
    \sym{$r$}{distance from the Earth's centre (the orbital radius)}{\si{m}}
    \sym{$r_\mathrm{E}$}{radius of the Earth, \num{6.371e6} (data sheet)}{\si{m}}
    \sym{$h$}{altitude: the height above the Earth's surface}{\si{m}}
  \end{symbols}
  \inwords{the distance from the centre is the Earth's radius plus the height above the ground.}
  \whentouse{whenever a question gives a height above the surface. Do it first, before any
    gravity formula. For another planet, use that planet's radius.}
  \whennottouse{when the question already gives the distance from the centre (the orbital radius).}
  \rearranged{$h = r - r_\mathrm{E}$ \quad (when a question asks for the altitude)}
  \tensecond{$h = \qty{300}{km} = \qty{0.300e6}{m}$, so
    $r = \qty{6.371e6}{m} + \qty{0.300e6}{m} = \qty{6.671e6}{m}$.}
  \derive{it is a definition (\cref{ch4:fig:altitude}). The line from the centre to the satellite
    crosses the surface. It has two parts: centre to surface ($r_\mathrm{E}$), then surface to
    satellite ($h$). Add the two parts.}
  \trap{change km to m first. Never put $h$ alone into $\dfrac{GM}{r^2}$, $\sqrt{\dfrac{GM}{r}}$
    or Kepler's third law.}
\end{formulacard}

\begin{formulacard}[notonsheet]{Gravitational field strength}
  \cardformula{g = \frac{GM}{r^2}}
  \begin{symbols}
    \sym{$g$}{gravitational field strength (free-fall acceleration)}{\si{N.kg^{-1}} = \si{m.s^{-2}}}
    \sym{$G$}{universal gravitational constant}{\si{N.m^2.kg^{-2}}}
    \sym{$M$}{mass of the planet (or star, or moon) that makes the field}{\si{kg}}
    \sym{$r$}{distance from the planet's \emph{centre} to the point}{\si{m}}
  \end{symbols}
  \inwords{$g$ is bigger for a heavier planet. It is smaller further from the centre: double $r$
    and $g$ is 4 times smaller. It does not depend on the object's mass.}
  \whentouse{$g$ at a planet's surface ($r$ = its radius) or above it ($r = r_\mathrm{E} + h$),
    and to compare planets.}
  \whennottouse{inside a planet.}
  \rearranged{$M = \dfrac{gr^2}{G}$ \qquad $r = \sqrt{\dfrac{GM}{g}}$}
  \tensecond{at $r = 2r_\mathrm{E}$, $g = \dfrac{9.86}{2^2} = \dfrac{9.86}{4} = \qty{2.47}{m.s^{-2}}$.}
  \derive{an object of mass $m$ feels $F = mg$ (definition of $g$) and $F = \dfrac{GMm}{r^2}$
    (on the sheet).
    \begin{align*}
      mg &= \frac{GMm}{r^2} && \text{both are the same force}\\
      g &= \frac{GM}{r^2} && \text{divide both sides by $m$; $m$ cancels}
    \end{align*}}
  \trap{$r$ is from the centre: $r = r_\mathrm{E} + h$, never the altitude $h$ alone. Change km to m.}
\end{formulacard}

\textbf{Two useful numbers.} Many questions need the product
$GM_\mathrm{E} = \num{6.67e-11} \times \num{6.0e24} = \qty{4.0e14}{m^3.s^{-2}}$
(\num{4.002e14} before rounding). At the surface, $r = r_\mathrm{E}$, so
$g = \dfrac{GM_\mathrm{E}}{r_\mathrm{E}^2} = \dfrac{\qty{4.002e14}{m^3.s^{-2}}}{(\qty{6.371e6}{m})^2}
= \qty{9.86}{m.s^{-2}}$.

This is a little more than the sheet's $g = \qty{9.8}{m.s^{-2}}$. The difference comes from
rounding $M_\mathrm{E}$ to \num{6.0e24}. The rule:
\begin{itemize}
  \item If a question gives $g$, or says ``use the data sheet'', use 9.8.
  \item If you compare a calculated $\dfrac{GM}{r^2}$ with the surface value, use 9.86. Then both
        numbers come from the same data.
\end{itemize}

\textbf{What affects $g$?} In the model $g = \dfrac{GM}{r^2}$ only two factors affect $g$: the
planet's mass ($g \propto M$) and the distance from its centre ($g \propto \dfrac{1}{r^2}$).
\Cref{ch4:fig:gr} shows the inverse-square fall. Doubling $r$ (an altitude of one Earth radius,
\qty{6371}{km}) makes $g$ four times smaller.

On the real Earth, $g$ also changes a little from place to place. It is about
\qty{9.78}{m.s^{-2}} at the equator and \qty{9.83}{m.s^{-2}} at the poles. There are two reasons.
The Earth bulges at the equator, so the surface there is further from the centre. The Earth also
spins, and this reduces the measured $g$ most at the equator. High mountains and dense rock
underground change $g$ slightly too.

\begin{figure}
  \centering
  \begin{tikzpicture}
    \begin{axis}[hsc, width=0.66\linewidth, height=0.30\linewidth,
        xmin=0, xmax=5.5, ymin=0, ymax=11,
        xlabel={distance from the Earth's centre, $r$ (in Earth radii)},
        ylabel={$g$ (\si{m.s^{-2}})}, xtick={0,1,2,3,4,5}, ytick={0,2,4,6,8,10}]
      \fill[hscmuted!15] (axis cs:0,0) rectangle (axis cs:1,11);
      \node[font=\footnotesize, text=hscmuted, align=center] at (axis cs:0.5,5.5) {inside\\Earth};
      \addplot[hscseriesA, line width=1.2pt, domain=1:5.5, samples=120, mark=none] {9.86/x^2};
      \addplot[hscseriesA, only marks, mark=*, mark size=2pt]
        coordinates {(1,9.86) (2,2.465) (3,1.096) (4,0.616)};
      \node[anchor=west, font=\small] at (axis cs:1.1,9.86) {9.86 at the surface};
      \node[anchor=south west, font=\small] at (axis cs:2.05,2.6) {2.47 ($\div 4$)};
      \node[anchor=south west, font=\small] at (axis cs:3.05,1.25) {1.10 ($\div 9$)};
      \node[anchor=south west, font=\small] at (axis cs:4.05,0.75) {0.62 ($\div 16$)};
    \end{axis}
  \end{tikzpicture}
  \caption{Field strength against distance from the Earth's centre (data-sheet values).}\label{ch4:fig:gr}
\end{figure}

\begin{workedexample}{Gravity on a weather satellite}\label{ch4:we:gsat}
A \qty{2500}{kg} weather satellite orbits \qty{850}{km} above the Earth's surface. Calculate
(a) the gravitational force on it and (b) the gravitational field strength there.
\given{$m = \qty{2500}{kg}$, \ $M_\mathrm{E} = \qty{6.0e24}{kg}$, \
  $G = \qty{6.67e-11}{N.m^2.kg^{-2}}$ (data sheet), \ add the altitude to the radius:
  $r = r_\mathrm{E} + h = \qty{6.371e6}{m} + \qty{0.850e6}{m} = \qty{7.221e6}{m}$}
\find{(a) the force $F$ \quad (b) the field strength $g$}
\formula{$F = \dfrac{GMm}{r^2}$, \quad $g = \dfrac{F}{m}$}
\substitute{\eqlines{GMm &= \qty{6.67e-11}{N.m^2.kg^{-2}} \times \qty{6.0e24}{kg} \times \qty{2500}{kg}\\
    &= \qty{1.0005e18}{N.m^2}\\
  F &= \frac{\qty{1.0005e18}{N.m^2}}{(\qty{7.221e6}{m})^2} = \qty{1.919e4}{N}\\
  g &= \frac{\qty{1.919e4}{N}}{\qty{2500}{kg}} = \qty{7.68}{N.kg^{-1}}}}
\answer{(a) \final{$F = \qty{1.92e4}{N}$} towards the Earth's centre. \ (b) \final{$g = \qty{7.68}{N.kg^{-1}}$},
  which is $7.68 \div 9.86 = 0.78$ (78\,\%) of the surface value.}
\checkstep{the other route:
  $g = \dfrac{GM}{r^2} = \dfrac{\qty{4.002e14}{m^3.s^{-2}}}{\qty{5.214e13}{m^2}} = \qty{7.68}{N.kg^{-1}}$,
  the same \cmark. (Here $r^2 = (\qty{7.221e6}{m})^2 = \qty{5.214e13}{m^2}$.) This is less than
  \qty{9.86}{N.kg^{-1}} at the surface. That makes sense, because the satellite is further from
  the centre.}
\end{workedexample}

\begin{toolkit}{Comparing two planets: divide one formula by the other}
Write the formula once for each planet. Then divide the two:
\[ \frac{g_\mathrm{X}}{g_\mathrm{E}} = \frac{GM_\mathrm{X}/r_\mathrm{X}^2}{GM_\mathrm{E}/r_\mathrm{E}^2}
   = \frac{M_\mathrm{X}}{M_\mathrm{E}} \times \frac{r_\mathrm{E}^2}{r_\mathrm{X}^2}
   = \frac{M_\mathrm{X}/M_\mathrm{E}}{(r_\mathrm{X}/r_\mathrm{E})^2} \]
$G$ is on the top and on the bottom, so it cancels. Dividing by a fraction is the same as
multiplying by it upside down. So you only need the two ratios, for example ``5.0 times the
mass'' and ``1.6 times the radius''.
\end{toolkit}

\begin{workedexample}{Weighing yourself on another planet}\label{ch4:we:planetx}
A rocky planet has 5.0 times the Earth's mass and 1.6 times the Earth's radius. (a) Find $g$ at
its surface. (b) Find the weight of a \qty{70}{kg} astronaut there. (c) An Earth bathroom scale
feels how hard you press on it. It shows that force divided by 9.8, and calls the result ``kg''.
What would an Earth scale show there?
\given{$\dfrac{M_\mathrm{X}}{M_\mathrm{E}} = 5.0$, \ $\dfrac{r_\mathrm{X}}{r_\mathrm{E}} = 1.6$, \
  $g_\mathrm{E} = \qty{9.8}{m.s^{-2}}$, \ $m = \qty{70}{kg}$}
\find{$g_\mathrm{X}$, the weight $W$, the scale reading}
\formula{ratio of $g = \dfrac{GM}{r^2}$ for the two planets ($G$ cancels; see the toolkit above):
  $\dfrac{g_\mathrm{X}}{g_\mathrm{E}} = \dfrac{M_\mathrm{X}/M_\mathrm{E}}{(r_\mathrm{X}/r_\mathrm{E})^2}$;
  \ $W = mg$}
\substitute{\eqlines{\frac{g_\mathrm{X}}{g_\mathrm{E}} &= \frac{5.0}{1.6^2} = \frac{5.0}{2.56} = 1.953\\
  g_\mathrm{X} &= 1.953 \times \qty{9.8}{m.s^{-2}} = \qty{19.1}{m.s^{-2}}\\
  W &= \qty{70}{kg} \times \qty{19.1}{m.s^{-2}} = \qty{1.34e3}{N}\\
  \text{reading} &= \qty{1.34e3}{N} \div \qty{9.8}{m.s^{-2}} = \qty{137}{kg}}}
\answer{(a) \final{$g_\mathrm{X} = \qty{19}{m.s^{-2}}$} \ (b) \final{$W = \qty{1.3e3}{N}$} \
  (c) about \final{\qty{1.4e2}{kg}} (2 significant figures). Her mass is still \qty{70}{kg}; only
  her weight has changed.}
\checkstep{the extra mass multiplies $g$ by 5.0; the bigger radius divides it by 2.56. The net
  effect is about $\times 2$, so the scale shows about twice \qty{70}{kg} \cmark.}
\end{workedexample}

\textbf{Fields from two bodies.} Field strength is a vector: it has a size and a direction. It
points towards the centre of the mass that makes it. When two bodies pull on the same point, add
their fields as vectors. On the line between two bodies the two fields point in opposite
directions, so you subtract their sizes.

\begin{workedexample}{A space probe between the Earth and the Moon}\label{ch4:we:earthmoon}
A space probe is on the line joining the centres of the Earth and the Moon
(\cref{ch4:fig:earthmoon}). It is \qty{3.00e8}{m} from the Earth's centre. The two centres are
\qty{3.84e8}{m} apart, and the Moon's mass is \qty{7.35e22}{kg}. Find the net gravitational
field strength at the probe.
\begin{figure}
  \centering
  \begin{tikzpicture}
    \fill[hscbrand!10, draw=hscink, line width=0.8pt] (0,0) circle (0.55);
    \node[diagramlabel, above, inner sep=4pt] at (0,0.55) {Earth};
    \fill[hscmuted!25, draw=hscink, line width=0.6pt] (9.0,0) circle (0.22);
    \node[diagramlabel, above, inner sep=4pt] at (9.0,0.22) {Moon};
    \fill[hscink] (0,0) circle (1.2pt);
    \fill[hscink] (9.0,0) circle (1.2pt);
    \draw[force] (7.0,0) -- (4.4,0) node[below, pos=0.5, text=hscforce, font=\small]
      {$g_\mathrm{E}$};
    \draw[force] (7.0,0) -- (7.42,0) node[below right, pos=0.6, text=hscforce, font=\small,
      inner sep=3pt] {$g_\mathrm{M}$};
    \fill[object] (7.0,0) circle (3.5pt);
    \node[diagramlabel, above, inner sep=5pt] at (7.0,0) {probe};
    \draw[guide] (0,0) -- (0,-1.0);
    \draw[guide] (7.0,0) -- (7.0,-1.0);
    \draw[guide] (9.0,0) -- (9.0,-1.0);
    \draw[hscmuted, line width=0.6pt, {Stealth[length=2mm]}-{Stealth[length=2mm]}]
      (0,-0.85) -- node[below, diagramlabel] {\qty{3.00e8}{m}} (7.0,-0.85);
    \draw[hscmuted, line width=0.6pt, {Stealth[length=2mm]}-{Stealth[length=2mm]}]
      (7.0,-0.85) -- node[below, diagramlabel] {\qty{0.84e8}{m}} (9.0,-0.85);
  \end{tikzpicture}
  \caption{A probe between the Earth and the Moon (the sizes of the bodies are not to scale). The
    two fields point in opposite directions. Arrow lengths show their sizes.}\label{ch4:fig:earthmoon}
\end{figure}
\given{$M_\mathrm{E} = \qty{6.0e24}{kg}$, \ $M_\mathrm{M} = \qty{7.35e22}{kg}$, \
  $G = \qty{6.67e-11}{N.m^2.kg^{-2}}$, \ distance to the Earth's centre $r_1 = \qty{3.00e8}{m}$, \
  distance to the Moon's centre $r_2 = \qty{3.84e8}{m} - \qty{3.00e8}{m} = \qty{0.84e8}{m}$}
\find{the net field strength $g_\mathrm{net}$ (size and direction)}
\formula{$g = \dfrac{GM}{r^2}$ for each body; opposite directions, so
  $g_\mathrm{net} = g_\mathrm{E} - g_\mathrm{M}$}
\substitute{\eqlines{g_\mathrm{E} &= \frac{\qty{6.67e-11}{N.m^2.kg^{-2}} \times \qty{6.0e24}{kg}}{(\qty{3.00e8}{m})^2}\\
    &= \qty{4.447e-3}{N.kg^{-1}} \quad\text{towards the Earth}\\
  g_\mathrm{M} &= \frac{\qty{6.67e-11}{N.m^2.kg^{-2}} \times \qty{7.35e22}{kg}}{(\qty{0.84e8}{m})^2}\\
    &= \qty{0.6948e-3}{N.kg^{-1}} \quad\text{towards the Moon}\\
  g_\mathrm{net} &= g_\mathrm{E} - g_\mathrm{M} \quad\text{(opposite directions, so subtract)}\\
    &= \qty{4.447e-3}{N.kg^{-1}} - \qty{0.6948e-3}{N.kg^{-1}}\\
    &= \qty{3.752e-3}{N.kg^{-1}}}}
\answer{\final{$g_\mathrm{net} = \qty{3.75e-3}{N.kg^{-1}}$ towards the Earth}}
\checkstep{the Earth has 82 times the Moon's mass, but the probe is much closer to the Moon. The
  Earth's field is still bigger here, so the net field points towards the Earth \cmark.}
\end{workedexample}

\begin{mistakes}
\begin{itemize}
  \item \xmark\ Using the altitude $h$ (or the surface-to-surface distance) as $r$.
        \cmark\ $r$ is from the centre: $r = r_\mathrm{E} + h$, in metres.
  \item \xmark\ Using $g = \qty{9.8}{m.s^{-2}}$ high above the Earth. \cmark\ Calculate
        $\dfrac{GM}{r^2}$ at that height.
  \item \xmark\ ``The Earth pulls the Moon harder than the Moon pulls the Earth.'' \cmark\ The
        forces are equal (third law). Only the accelerations differ.
  \item \xmark\ ``There is no gravity in space.'' \cmark\ At \qty{420}{km}, $g$ is still
        \qty{8.7}{m.s^{-2}}.
\end{itemize}
\end{mistakes}

\begin{examtip}
Calculations are usually 2 marks: one for substituting with $r$ in metres from the centre, one
for the answer with a unit. When a planet's mass and radius are given as multiples of the
Earth's, use the ratio form. Multiple-choice questions often change two things at once. Write
$F \propto \dfrac{Mm}{r^2}$ (or $g \propto \dfrac{M}{r^2}$). Then change one factor at a time.
$g = \dfrac{GM}{r^2}$ is \emph{not} on the sheet, so derive it in one line from $mg = \dfrac{GMm}{r^2}$.
Asked why $g$ is different in different places on the Earth? Give altitude, latitude (the
Earth's shape and spin) and the density of the rock below.
\end{examtip}

\begin{quickcheck}
  \item Two asteroids attract each other with a force $F$. The distance between their centres
        is tripled and the mass of one asteroid is doubled. Find the new force in terms of $F$.
  \item The Earth pulls a \qty{0.50}{kg} ball down with a force of \qty{4.9}{N}. (a) What force
        does the ball exert on the Earth? (b) Explain why the Earth does not seem to move.
  \item The Moon has a mass of \qty{7.35e22}{kg} and a radius of \qty{1.74e6}{m}. Calculate $g$
        at its surface.
  \item At what altitude above the Earth's surface is $g$ one ninth of its surface value?
  \item Use the data in \cref{ch4:we:earthmoon}. How far from the Earth's centre is the net field
        of the Earth and the Moon zero?
\end{quickcheck}

%% ==================================================================
\section{Satellites in circular orbits}\label{ch4:sec:orbits}
\covers{M5.IQ3.02, M5.IQ3.04}

\begin{bigidea}
A satellite is always falling towards the Earth. It moves sideways so fast that the ground
curves away beneath it, so it never lands. Gravity provides the centripetal force.
\end{bigidea}

\begin{keywords}
  \kw{Satellite}{an object that orbits a bigger body. The Moon is a natural satellite.}
  \kw{Orbital radius ($r$)}{the distance from the centre of the planet to the satellite.}
  \kw{Orbital speed ($v$)}{the speed a satellite needs to stay in a circular orbit.}
  \kw{Orbital period ($T$)}{the time for one full orbit.}
  \kw{Free fall}{motion when gravity is the only force acting.}
\end{keywords}

\begin{refresher}{Circular motion (\cref{ch2:chap})}
In uniform circular motion, $v = \dfrac{2\pi r}{T}$, the acceleration $a_\mathrm{c} = \dfrac{v^2}{r}$
points to the centre, and the net force is $F_\mathrm{c} = \dfrac{mv^2}{r}$. $F_\mathrm{c}$ is not a
new force: a real force (tension, friction, gravity) does this job.
\end{refresher}

\begin{sidebyside}[0.5]
  \centering
  \includegraphics{ch4_cannonball.pdf}
  \captionof{figure}{Newton's cannonball (the mountain is drawn far too tall, so the paths can be
    seen). C is a circular orbit: gravity $F$ is at \ang{90} to $v$.}\label{ch4:fig:cannon}
\nextside
  Newton imagined a cannon on a very tall mountain, with no air (\cref{ch4:fig:cannon}).
  \begin{itemize}
    \item A slow ball falls to the ground (A). A faster one lands further away (B).
    \item At about \qty{7.9}{km.s^{-1}} near the surface, the ground curves away as fast as the
          ball falls. It never lands: it is in \term{orbit} (C).
    \item Faster still, the orbit is an oval (D). Above \qty{11.2}{km.s^{-1}} the ball escapes
          (\cref{ch5:chap}).
  \end{itemize}
  An orbit is projectile motion that never ends. The satellite is in free fall.
\end{sidebyside}

Gravity points to the Earth's centre, exactly where a centripetal force must point. So
\emph{gravity provides the centripetal force}: $\dfrac{GMm}{r^2} = \dfrac{mv^2}{r}$.

\begin{formulacard}[notonsheet]{Orbital speed (circular orbit)}
  \cardformula{v = \sqrt{\frac{GM}{r}}}
  \begin{symbols}
    \sym{$v$}{orbital speed of the satellite}{\si{m.s^{-1}}}
    \sym{$G$}{universal gravitational constant}{\si{N.m^2.kg^{-2}}}
    \sym{$M$}{mass of the central body (the planet, \emph{not} the satellite)}{\si{kg}}
    \sym{$r$}{orbital radius, from the centre of the central body}{\si{m}}
  \end{symbols}
  \inwords{a satellite closer to the planet must move faster.}
  \whentouse{any circular orbit around a planet, moon or star.}
  \whennottouse{an elliptical orbit (the speed changes), or a rocket during launch.}
  \rearranged{$r = \dfrac{GM}{v^2}$ \qquad $M = \dfrac{v^2 r}{G}$}
  \tensecond{around the Earth at $r = \qty{7.0e6}{m}$:
    $v = \sqrt{\dfrac{\qty{4.0e14}{m^3.s^{-2}}}{\qty{7.0e6}{m}}} = \qty{7.6e3}{m.s^{-1}}$.}
  \derive{set $F = F_\mathrm{c}$ (both on the sheet).
    \begin{align*}
      \frac{GMm}{r^2} &= \frac{mv^2}{r} && \text{gravity provides the centripetal force}\\
      \frac{GM}{r^2} &= \frac{v^2}{r} && \text{divide both sides by $m$}\\
      \frac{GM}{r} &= v^2 && \text{multiply both sides by $r$}\\
      v &= \sqrt{\frac{GM}{r}} && \text{take the square root}
    \end{align*}}
  \trap{$M$ is the planet's mass, not the satellite's. $r$ is not the altitude.}
\end{formulacard}

The derivation shows four facts.
\begin{enumerate}
  \item The satellite's mass cancels. So every object at the same height needs the same speed.
        A spanner dropped by an astronaut orbits beside the station.
  \item A bigger orbit is slower: $v \propto \dfrac{1}{\sqrt r}$. For example, 9 times the radius
        gives one third of the speed, because $\sqrt 9 = 3$.
  \item Put $v^2 = \dfrac{GM}{r}$ into $a_\mathrm{c} = \dfrac{v^2}{r}$:
        $a_\mathrm{c} = \dfrac{GM}{r} \div r = \dfrac{GM}{r^2}$. This is the field strength $g$ at
        the orbit.
  \item The velocity is along the tangent to the orbit, at \ang{90} to the radius. If a question
        asks for the orbital \emph{velocity}, give the speed and say ``along the tangent to the
        orbit''.
\end{enumerate}

\term{Apparent weightlessness} is the feeling of having no weight during free fall. Astronauts on
the International Space Station (ISS, about \qty{420}{km} up) float. This is not because gravity is
missing: $g$ there is \qty{8.7}{m.s^{-2}}. We feel our weight because the floor pushes up on us.
This push is the \term{normal force}. In orbit, the astronaut and the station fall together with
the same acceleration. So the floor does not push on the astronaut: the normal force is zero, and
the astronaut feels weightless.

\subsection{The period and Kepler's third law}
In one orbit a satellite travels $2\pi r$ in time $T$, so $v = \dfrac{2\pi r}{T}$ (on the
sheet). Joining this with the orbital speed links $r$ and $T$. This is Kepler's third law in the
form Newton found (\cref{ch4:sec:kepler}). Johannes Kepler was an astronomer. In 1619 he found
this rule from planet data.

\begin{toolkit}{Cube roots, the power 3/2 and seconds}
\begin{itemize}
  \item Cube root: use the $\sqrt[3]{\phantom{x}}$ key. $\sqrt[3]{\num{8.0e21}} = \num{2.0e7}$,
        because $(\num{2.0e7})^3 = \num{8.0e21}$.
  \item $x^{3/2} = \sqrt{x^3}$: \ $9^{3/2} = \sqrt{729} = 27$. So 9 times the radius gives 27 times
        the period.
  \item \qty{1}{min} = \qty{60}{s}, \ \qty{1}{h} = \qty{3600}{s}, \ 1 day = \qty{86400}{s}.
        Also, $4\pi^2 = 39.48$.
\end{itemize}
\end{toolkit}

\begin{formulacard}{Kepler's third law (Newton's form)}
  \cardformula{\frac{r^3}{T^2} = \frac{GM}{4\pi^2}}
  \begin{symbols}
    \sym{$r$}{orbital radius (for an ellipse: the semi-major axis, \cref{ch4:sec:kepler})}{\si{m}}
    \sym{$T$}{period: the time for one orbit}{\si{s}}
    \sym{$G$}{universal gravitational constant}{\si{N.m^2.kg^{-2}}}
    \sym{$M$}{mass of the central body}{\si{kg}}
  \end{symbols}
  \inwords{$\dfrac{r^3}{T^2}$ is the same for everything that orbits the same body.}
  \whentouse{to find $T$ from $r$, $r$ from $T$, or the central mass $M$ from $r$ and $T$.
    For two objects orbiting the same body: $\dfrac{r_1^3}{T_1^2} = \dfrac{r_2^3}{T_2^2}$.}
  \whennottouse{to compare objects that orbit different central bodies. For example, Titan
    (orbits Saturn) and Venus (orbits the Sun) have different values of $\dfrac{r^3}{T^2}$.}
  \rearranged{$T = 2\pi\sqrt{\dfrac{r^3}{GM}}$ \qquad $r = \sqrt[3]{\dfrac{GMT^2}{4\pi^2}}$ \qquad
    $M = \dfrac{4\pi^2 r^3}{GT^2}$ \quad (steps in the toolkit below)}
  \tensecond{at $r = \qty{7.0e6}{m}$ around the Earth,
    $T = 2\pi\sqrt{\dfrac{(\qty{7.0e6}{m})^3}{\qty{4.0e14}{m^3.s^{-2}}}} = \qty{5.8e3}{s}
    = \qty{97}{min}$.}
  \derive{$v^2 = \dfrac{GM}{r}$ (orbital speed) and $v = \dfrac{2\pi r}{T}$ (on the sheet).
    \begin{align*}
      v^2 &= \frac{4\pi^2 r^2}{T^2} && \text{square both sides of $v = \frac{2\pi r}{T}$}\\
      \frac{4\pi^2 r^2}{T^2} &= \frac{GM}{r} && \text{both expressions equal $v^2$}\\
      \frac{4\pi^2 r^3}{T^2} &= GM && \text{multiply both sides by $r$}\\
      \frac{r^3}{T^2} &= \frac{GM}{4\pi^2} && \text{divide both sides by $4\pi^2$}
    \end{align*}}
  \trap{$T$ must be in seconds. Use a cube root (not a square root) to find $r$, then subtract
    $r_\mathrm{E}$ for the altitude.}
\end{formulacard}

\begin{toolkit}{Rearranging Kepler's third law}
\begin{align*}
  \frac{r^3}{T^2} &= \frac{GM}{4\pi^2} && \text{start (on the sheet)}\\
  r^3 &= \frac{GMT^2}{4\pi^2} && \text{multiply both sides by $T^2$}\\
  4\pi^2 r^3 &= GMT^2 && \text{multiply both sides by $4\pi^2$}\\
  T^2 &= \frac{4\pi^2 r^3}{GM} && \text{divide both sides by $GM$}\\
  T &= 2\pi\sqrt{\frac{r^3}{GM}} && \text{square root of both sides; $\sqrt{4\pi^2} = 2\pi$}
\end{align*}
To find $r$: take the cube root of both sides of the second line. To find $M$: divide both sides
of the third line by $GT^2$.
\end{toolkit}

\begin{workedexample}{A satellite 600 km up}\label{ch4:we:sat600}
A satellite moves in a circular orbit \qty{600}{km} above the Earth's surface. Calculate (a) its
orbital speed and (b) its period in minutes.
\given{$r = \qty{6.371e6}{m} + \qty{0.600e6}{m} = \qty{6.971e6}{m}$, \
  $GM_\mathrm{E} = \num{6.67e-11} \times \num{6.0e24} = \qty{4.002e14}{m^3.s^{-2}}$}
\find{(a) $v$ \quad (b) $T$}
\formula{$v = \sqrt{\dfrac{GM}{r}}$, \quad $T = 2\pi\sqrt{\dfrac{r^3}{GM}}$ (Kepler's third law)}
\substitute{\eqlines{v &= \sqrt{\frac{\qty{4.002e14}{m^3.s^{-2}}}{\qty{6.971e6}{m}}}
    = \sqrt{\qty{5.741e7}{m^2.s^{-2}}}\\
    &= \qty{7.577e3}{m.s^{-1}}\\
  T &= 2\pi\sqrt{\frac{(\qty{6.971e6}{m})^3}{\qty{4.002e14}{m^3.s^{-2}}}} = 2\pi\sqrt{\qty{8.465e5}{s^2}}\\
    &= 2\pi \times \qty{920.0}{s} = \qty{5.78e3}{s}\\
  T &= \qty{5.78e3}{s} \div 60 = \qty{96.3}{min}}}
\answer{(a) \final{$v = \qty{7.58e3}{m.s^{-1}}$} \ (b) \final{$T = \qty{96.3}{min}$}}
\checkstep{$v = \dfrac{2\pi r}{T} = \dfrac{2\pi \times \num{6.971e6}}{\num{5.78e3}} = \qty{7.58e3}{m.s^{-1}}$,
  the same speed \cmark. Also $a_\mathrm{c} = \dfrac{v^2}{r} = \qty{8.24}{m.s^{-2}} = \dfrac{GM}{r^2}$,
  the local $g$ \cmark.}
\end{workedexample}

\begin{workedexample}{Finding the mass of Saturn}\label{ch4:we:saturn}
Titan, Saturn's largest moon, has a near-circular orbit of radius \qty{1.22e9}{m} and a period of
15.9 days. Calculate the mass of Saturn.
\given{$r = \qty{1.22e9}{m}$, \ $T = 15.9 \times \qty{86400}{s} = \qty{1.374e6}{s}$, \
  $G = \qty{6.67e-11}{N.m^2.kg^{-2}}$ (data sheet)}
\find{the mass $M$ of Saturn (the central body)}
\formula{$\dfrac{r^3}{T^2} = \dfrac{GM}{4\pi^2}$, \ rearranged: $M = \dfrac{4\pi^2 r^3}{GT^2}$}
\substitute{\eqlines{M &= \frac{39.48 \times (\qty{1.22e9}{m})^3}{\qty{6.67e-11}{N.m^2.kg^{-2}} \times (\qty{1.374e6}{s})^2}\\
  &= \frac{\qty{7.169e28}{m^3}}{\qty{125.9}{N.m^2.kg^{-2}.s^2}} = \qty{5.69e26}{kg}}}
\answer{\final{$M = \qty{5.69e26}{kg}$} (3 significant figures). Keep all the calculator digits
  until the end, then round.}
\checkstep{this is $\num{5.69e26} \div \num{6.0e24} \approx 95$ Earth masses, sensible for a giant
  planet \cmark. (The accepted value is \qty{5.68e26}{kg}.)}
\end{workedexample}

\begin{mistakes}
\begin{itemize}
  \item \xmark\ Putting the satellite's mass in for $M$. \cmark\ $M$ is the central body's mass.
  \item \xmark\ Using minutes, hours or days for $T$ with $G$. \cmark\ Change $T$ to seconds.
  \item \xmark\ Taking a square root to find $r$ from $r^3$. \cmark\ Take a cube root.
  \item \xmark\ ``Astronauts float because there is no gravity.'' \cmark\ They are in free fall
        with their spacecraft.
\end{itemize}
\end{mistakes}

\begin{examtip}
``Derive $v = \sqrt{\dfrac{GM}{r}}$'' or ``show that $\dfrac{r^3}{T^2} = \dfrac{GM}{4\pi^2}$''
(2--3 marks):
\begin{enumerate}
  \item Write ``gravity provides the centripetal force''.
  \item Write $\dfrac{GMm}{r^2} = \dfrac{mv^2}{r}$.
  \item Show that $m$ cancels. Write one step per line.
\end{enumerate}
``Explain why astronauts feel weightless'' (2--3 marks). Give $g$ in orbit (about
\qty{8.7}{m.s^{-2}}). Say that the astronaut and the craft are both in free fall with the same
acceleration. Say that the normal force is zero.
\end{examtip}

\begin{quickcheck}
  \item Calculate the orbital speed of a satellite \qty{2000}{km} above the Earth's surface.
  \item Satellite B orbits the Earth at 4 times the orbital radius of satellite A. Compare their
        speeds and their periods.
  \item The Moon orbits the Earth at a radius of \qty{3.84e8}{m}. Calculate its period in days.
  \item An astronaut on the ISS feels weightless, even though $g = \qty{8.7}{m.s^{-2}}$ there.
        Explain why.
\end{quickcheck}

%% ==================================================================
\section{Kepler's laws of planetary motion}\label{ch4:sec:kepler}
\covers{M5.IQ3.04}

\begin{bigidea}
Kepler found three rules for how planets move by studying careful measurements. Newton's law of
gravitation explains all three. In an oval orbit, gravity speeds a planet up as it moves closer
to the Sun and slows it down as it moves away.
\end{bigidea}

\begin{keywords}
  \kw{Ellipse}{a stretched circle (an oval).}
  \kw{Focus (plural: foci)}{one of two fixed points inside an ellipse, on its long axis. From any
    point on the ellipse, the distances to the two foci add up to the same total.}
  \kw{Semi-major axis ($a$)}{half of the longest diameter of an ellipse.}
  \kw{Eccentricity ($e$)}{how stretched an ellipse is: 0 is a circle, close to 1 is very flat.}
  \kw{Perihelion and aphelion}{the closest and farthest points of an orbit around the Sun
    (around the Earth: perigee and apogee).}
\end{keywords}

In the late 1500s Tycho Brahe measured the positions of the planets very carefully, without a
telescope. Johannes Kepler used Tycho's data on Mars to find his first two laws (1609). In 1619
he found the third law by comparing the periods and distances of all the known planets. They are
\term{empirical laws}: rules found from data, which describe the motion but do not explain it. In
1687 Newton showed that a force with $F \propto \dfrac{1}{r^2}$ explains all three.

\subsection{First law: the law of ellipses}
\begin{sidebyside}[0.56]
  \centering
  \begin{tikzpicture}[scale=0.74]
    \draw[hscink, line width=0.9pt] (0,0) ellipse (3 and 2.598);
    \draw[guide] (-3,0) -- (3,0);
    \fill[hscink] (0,0) circle (1.6pt);
    \node[diagramlabel, below, inner sep=2pt] at (0,0) {centre};
    \draw[hscink, line width=0.6pt] (-1.5,0) circle (2.2pt);
    \node[diagramlabel, above, align=center, inner sep=4pt] at (-1.5,0.12) {empty\\focus};
    \fill[hscray, draw=hscink, line width=0.5pt] (1.5,0) circle (5pt);
    \node[diagramlabel, below, inner sep=3pt] at (1.5,-0.15) {Sun};
    \draw[hscmuted, line width=0.6pt, {Stealth[length=2mm]}-{Stealth[length=2mm]}]
      (0,-1.35) -- node[below, diagramlabel] {$a$} (3,-1.35);
    \fill[hscvel] (3,0) circle (2.2pt);
    \node[diagramlabel, right] at (3.05,0) {perihelion};
    \fill[hscvel] (-3,0) circle (2.2pt);
    \node[diagramlabel, left] at (-3.05,0) {aphelion};
    \fill[object] (-1.5,2.25) circle (4.5pt);
    \node[diagramlabel, above left] at (-1.55,2.3) {planet};
  \end{tikzpicture}
  \captionof{figure}{An ellipse with $e = 0.5$, to scale. $a$ is the semi-major axis.}\label{ch4:fig:ellipse}
\nextside
  \emph{Each planet moves in an ellipse, with the Sun at one focus} (\cref{ch4:fig:ellipse}). The
  Sun is not at the centre, and nothing sits at the other focus.

  \smallskip
  Most planet orbits are nearly circles (the Earth's $e$ is only 0.017), so HSC calculations treat
  them as circles. Comets have very stretched orbits. For an ellipse, the $r$ in Kepler's third
  law is the semi-major axis: $a = \dfrac{r_\mathrm{p} + r_\mathrm{a}}{2}$ (perihelion and
  aphelion distances from the Sun).
\end{sidebyside}

\subsection{Second law: the law of equal areas}
\begin{sidebyside}[0.56]
  \centering
  \includegraphics{ch4_equalareas.pdf}
  \captionof{figure}{Equal times, equal areas: $\Delta t_1 = \Delta t_2$, so $A_1 = A_2$ (to scale).}
  \label{ch4:fig:areas}
\nextside
  \emph{A line from the Sun to a planet sweeps out equal areas in equal times}
  (\cref{ch4:fig:areas}).

  \smallskip
  Near the Sun the line is short. To sweep the same area, the planet must cover a longer arc in
  the same time. So a planet moves \textbf{fastest at perihelion} and \textbf{slowest at
  aphelion}.
\end{sidebyside}

\begin{refresher}{Work and energy (\cref{ch3:chap})}
\term{Work} is energy given to an object by a force: $W = Fs\cos\theta$, where $\theta$ is the
angle between the force and the motion. Positive work speeds an object up. Negative work slows it
down. \term{Kinetic energy} $K = \tfrac12 mv^2$ is the energy of motion.
\end{refresher}

\begin{sidebyside}[0.56]
  \centering
  \includegraphics{ch4_ellipseforces.pdf}
  \captionof{figure}{Velocity (blue): the arrow length shows the speed. Gravity (red): these
    arrows show only the direction. At aphelion gravity is really 9 times weaker than at
    perihelion. Dashed red: the part of gravity along the motion.}\label{ch4:fig:forces}
\nextside
  \textbf{With forces} (\cref{ch4:fig:forces}). A force does work only when part of it points
  along the motion ($W = Fs\cos\theta$).
  \begin{itemize}
    \item \textbf{Circular orbit:} gravity is always at \ang{90} to $v$, so it does no work
          ($\cos 90^\circ = 0$). The speed is constant.
    \item \textbf{Moving closer:} part of gravity points along the motion and does work, so the
          planet speeds up.
    \item \textbf{Moving away:} part of gravity points against the motion, so it slows down.
    \item At perihelion and aphelion gravity is at \ang{90} to $v$: the speed is at its largest
          or smallest.
  \end{itemize}
\end{sidebyside}

\textbf{With energy.} Only gravity acts on the planet. No other force adds or removes energy. So
the planet's \term{total energy} (kinetic energy $K$ plus gravitational potential energy $U$)
stays constant. The sheet gives $U = -\dfrac{GMm}{r}$. (\Cref{ch5:chap} explains the minus sign.)
A smaller $r$ makes $\dfrac{GMm}{r}$ bigger, so $U$ is more negative. For example, $-10$ is lower
than $-5$. So $U$ is lowest when the planet is closest to the Sun. The total stays the same, so
$K$ is largest there, and the planet is fastest. In a circular orbit $r$ is constant, so $U$, $K$
and the speed are all constant.

% layout QA: short title on the continuation part, so the title bar stays on one line
\begin{formulacard}[notonsheet, title after break={\hscboxtitle{\faIcon{square-root-alt}}{Formula card}{Kepler's second law}\hspace{0.5em}{\footnotesize\mdseries\itshape(continued)}}]{Kepler's second law at the closest and farthest points}
  \cardformula{r_\mathrm{p} v_\mathrm{p} = r_\mathrm{a} v_\mathrm{a}}
  \begin{symbols}
    \sym{$r_\mathrm{p}$, $r_\mathrm{a}$}{distance from the centre of the central body (Sun,
      Earth, \ldots) at the closest point and at the farthest point}{\si{m}}
    \sym{$v_\mathrm{p}$, $v_\mathrm{a}$}{speed at the closest point (perihelion or perigee) and
      at the farthest point (aphelion or apogee)}{\si{m.s^{-1}}}
  \end{symbols}
  \inwords{at the two ends of an orbit, distance $\times$ speed is the same.}
  \whentouse{only at perihelion and aphelion (or perigee and apogee) of \emph{one} orbit.}
  \whennottouse{at other points of the orbit, or to compare two different planets.}
  \rearranged{$v_\mathrm{a} = \dfrac{r_\mathrm{p} v_\mathrm{p}}{r_\mathrm{a}}$ \qquad
    $r_\mathrm{a} = \dfrac{r_\mathrm{p} v_\mathrm{p}}{v_\mathrm{a}}$}
  \tensecond{twice as far away means half as fast: $r_\mathrm{a} = 2r_\mathrm{p}$ gives
    $v_\mathrm{a} = \tfrac12 v_\mathrm{p}$.}
  \derive{at perihelion and aphelion, $v$ is at \ang{90} to the line from the Sun. In a short time
    $\Delta t$, the planet moves a distance $v\Delta t$. The line from the Sun sweeps out a thin
    triangle. Its base is $r$ and its height is $v\Delta t$. So its area is
    $\tfrac12 \times \text{base} \times \text{height} = \tfrac12 r v \Delta t$.
    \begin{align*}
      \tfrac12\, r_\mathrm{p} v_\mathrm{p}\Delta t &= \tfrac12\, r_\mathrm{a} v_\mathrm{a}\Delta t
        && \text{equal areas in the same time}\\
      r_\mathrm{p} v_\mathrm{p} &= r_\mathrm{a} v_\mathrm{a} && \text{cancel $\tfrac12$ and $\Delta t$}
    \end{align*}}
  \trap{measure both distances from the Sun (the focus), not from the centre of the ellipse.}
\end{formulacard}

\subsection{Third law: the law of periods}
\emph{The square of the period is proportional to the cube of the orbital radius (semi-major
axis)}: $T^2 \propto r^3$. Newton's form (\cref{ch4:sec:orbits}) is
$\dfrac{r^3}{T^2} = \dfrac{GM}{4\pi^2}$, so the constant depends only on the central mass. All the
Sun's planets share one value; all the Earth's satellites share another. The ratio form
$\dfrac{r_1^3}{T_1^2} = \dfrac{r_2^3}{T_2^2}$ works in any units, if both orbits use the same
units. For example, use years and \term{astronomical units} (AU). One AU is the average
Earth--Sun distance, about \qty{1.50e11}{m}.

\begin{workedexample}{The speed of a comet}\label{ch4:we:comet}
A comet orbits the Sun. At perihelion it is \qty{1.2e11}{m} from the Sun's centre and moves at
\qty{4.29e4}{m.s^{-1}}. At aphelion it is \qty{6.0e11}{m} from the Sun's centre. Calculate its
speed at aphelion.
\given{$r_\mathrm{p} = \qty{1.2e11}{m}$, \ $v_\mathrm{p} = \qty{4.29e4}{m.s^{-1}}$, \
  $r_\mathrm{a} = \qty{6.0e11}{m}$}
\find{the speed $v_\mathrm{a}$ at aphelion}
\formula{$r_\mathrm{p} v_\mathrm{p} = r_\mathrm{a} v_\mathrm{a}$ (the two ends of one orbit), so
  $v_\mathrm{a} = \dfrac{r_\mathrm{p} v_\mathrm{p}}{r_\mathrm{a}}$}
\substitute{$v_\mathrm{a} = \dfrac{\qty{1.2e11}{m} \times \qty{4.29e4}{m.s^{-1}}}{\qty{6.0e11}{m}}
  = \qty{8.58e3}{m.s^{-1}}$}
\answer{\final{$v_\mathrm{a} = \qty{8.6e3}{m.s^{-1}}$} (2 significant figures).}
\checkstep{5 times further away, so 5 times slower \cmark. The comet is slower far from the Sun,
  as Kepler's second law says.}
\end{workedexample}

\begin{workedexample}{The period of the comet}\label{ch4:we:cometT}
Find the period of the comet in \cref{ch4:we:comet}. The Earth orbits the Sun at a radius of
\qty{1.50e11}{m} with a period of 1.00 year.
\given{comet: $r_\mathrm{p} = \qty{1.2e11}{m}$, $r_\mathrm{a} = \qty{6.0e11}{m}$; \
  Earth: $r_1 = \qty{1.50e11}{m}$, $T_1 = \qty{1.00}{\yr}$}
\find{the comet's period $T_2$, in years}
\formula{\eqlines{a &= \frac{r_\mathrm{p} + r_\mathrm{a}}{2} \quad\text{(semi-major axis)}\\
  \frac{r_1^3}{T_1^2} &= \frac{r_2^3}{T_2^2} \quad\text{(same central body: the Sun)}\\
  T_2^2\, r_1^3 &= T_1^2\, r_2^3 \quad\text{(multiply both sides by $T_1^2 T_2^2$)}\\
  T_2^2 &= T_1^2 \times \frac{r_2^3}{r_1^3} = T_1^2\left(\frac{r_2}{r_1}\right)^3
    \quad\text{(divide both sides by $r_1^3$)}}}
\substitute{\eqlines{r_2 = a &= \frac{\qty{1.2e11}{m} + \qty{6.0e11}{m}}{2} = \qty{3.6e11}{m}\\
  T_2^2 &= (\qty{1.00}{\yr})^2 \times \left(\frac{\qty{3.6e11}{m}}{\qty{1.50e11}{m}}\right)^3\\
    &= (\qty{1.00}{\yr})^2 \times 2.4^3 = \qty{13.82}{\yr^2}
    \quad\text{(because $3.6 \div 1.50 = 2.4$)}\\
  T_2 &= \sqrt{\qty{13.82}{\yr^2}} = \qty{3.72}{\yr} \quad\text{(square root)}}}
\answer{\final{$T_2 = \qty{3.7}{\yr}$} (2 significant figures).}
\checkstep{the comet's semi-major axis is 2.4 times the Earth's. Mars (1.5 times) takes
  \qty{1.9}{\yr}. Jupiter (5.2 times) takes \qty{12}{\yr}. The comet's \qty{3.7}{\yr} lies between
  them \cmark.}
\end{workedexample}

\begin{mistakes}
\begin{itemize}
  \item \xmark\ Putting the Sun at the centre of the ellipse. \cmark\ The Sun is at one focus.
  \item \xmark\ ``Gravity is stronger near the Sun, so it pushes the planet faster.'' \cmark\ Give
        the reason: part of gravity acts along the motion (work), or $U$ changes into $K$.
  \item \xmark\ Using $r_\mathrm{p}v_\mathrm{p} = r_\mathrm{a}v_\mathrm{a}$ at any two points.
        \cmark\ Only at perihelion and aphelion.
  \item \xmark\ Using the perihelion distance as $r$ in the third law. \cmark\ Use the
        semi-major axis.
\end{itemize}
\end{mistakes}

\begin{examtip}
``Explain why a planet moves fastest at perihelion'' (3--4 marks). Link three ideas:
\begin{enumerate}
  \item Equal areas in equal times. Near the Sun the line is short, so the planet must cover a
        longer arc in the same time.
  \item Total energy is constant. $U$ is lowest at perihelion, so $K$ is highest there.
  \item On the way in, part of gravity points along the motion. It does work and speeds the
        planet up.
\end{enumerate}
``Relate Kepler's third law to Newton's law'' needs the derivation from $F = F_\mathrm{c}$.
\end{examtip}

\begin{quickcheck}
  \item A satellite's perigee and apogee are \qty{7.0e6}{m} and \qty{1.4e7}{m} from the Earth's
        centre. Its speed at perigee is \qty{8.73e3}{m.s^{-1}}. Find its speed at apogee.
  \item The dwarf planet Ceres orbits the Sun with a semi-major axis 2.77 times the Earth's. Find
        its period in Earth years.
  \item Use forces to explain why a planet in a circular orbit moves at a constant speed, but a
        planet in an elliptical orbit does not.
  \item Use energy to explain why a comet moves most slowly at aphelion.
\end{quickcheck}

%% ==================================================================
\section{Low Earth orbits and geostationary orbits}\label{ch4:sec:leogeo}
\covers{M5.IQ3.03}

\begin{bigidea}
The height of an orbit fixes its speed and period. Low satellites circle the Earth in about
90 minutes, while geostationary satellites take one day and seem to hang still. Each height suits
different jobs.
\end{bigidea}

\begin{keywords}
  \kw{Low Earth orbit (LEO)}{an orbit about \qtyrange{200}{2000}{km} above the Earth's surface.}
  \kw{Geostationary orbit}{an orbit in which a satellite stays above one point on the equator.}
  \kw{Geosynchronous orbit}{any orbit with a period of one day. It may be tilted.}
  \kw{Orbital decay}{the slow loss of height caused by drag from the thin upper atmosphere.}
\end{keywords}

\textbf{Low Earth orbit.} At \qtyrange{200}{2000}{km}, a satellite moves at
\qtyrange{6.9}{7.8}{km.s^{-1}} and orbits in \qtyrange{88}{127}{min}. Thin air still drags on it,
so it slowly sinks (\term{orbital decay}, \cref{ch5:chap}). A \emph{polar} LEO passes over both
poles. The Earth turns beneath it, so it passes over every part of the Earth. A wide-view weather
camera sees the whole Earth within a day. A narrow, detailed camera takes a few days or more.

\textbf{Geostationary orbit.} A \term{geostationary} satellite needs all three conditions:
\begin{enumerate}
  \item a period of one day (strictly 23\,h\,56\,min, the time for the Earth to spin once);
  \item a circular orbit in the plane of the equator;
  \item motion from west to east, the same way as the Earth spins.
\end{enumerate}
Kepler's third law then allows only one radius, \qty{4.23e7}{m} (\cref{ch4:we:geo}). A
\term{geosynchronous} orbit has the one-day period but may be tilted. Such a satellite drifts
north and south in the sky each day, so it is not geostationary.

\begin{sidebyside}[0.47]
  \centering
  \begin{tikzpicture}[scale=0.86]
    % radii: Earth radius = 0.55 cm; ISS 1.066, GPS 4.18, GEO 6.64 Earth radii
    \draw[hscseriesB, line width=0.9pt] (0,0) circle (3.652);
    \draw[hscseriesC, line width=0.8pt, densely dotted] (0,0) circle (2.300);
    \fill[hscbrand!15, draw=hscink, line width=0.7pt] (0,0) circle (0.55);
    \draw[hscseriesA, line width=0.8pt] (0,0) circle (0.586);
    \draw[->, hscink, line width=0.6pt] (60:0.3) arc (60:300:0.3);
    \fill[object] (45:3.652) circle (3.5pt);
    \draw[hscseriesB, ->, line width=1pt] (52:3.652) arc (52:72:3.652);
    \node[diagramlabel, text=hscseriesB, font=\small\bfseries] at (0,3.25) {GEO};
    \node[diagramlabel, text=hscseriesC, font=\small\bfseries] at (0,1.95) {GPS};
    \draw[hscseriesA, line width=0.5pt] (0.62,0) -- (1.05,0);
    \node[diagramlabel, text=hscseriesA, font=\small\bfseries, anchor=west, inner sep=1pt] at (1.05,0) {LEO};
  \end{tikzpicture}
  \captionof{figure}{Orbit sizes to scale, seen from above the North Pole (the tilts of the ISS
    and GPS orbits are ignored). The Earth spins anticlockwise, and the satellites also move
    anticlockwise.}\label{ch4:fig:orbitscale}
\nextside
  \Cref{ch4:fig:orbitscale} shows three orbits to scale:
  \begin{itemize}
    \item \textcolor{hscseriesA}{\bfseries LEO} (ISS): \qty{420}{km} up, \qty{93}{min}. It almost
          touches the Earth.
    \item \textcolor{hscseriesC}{\bfseries GPS}: about \qty{20200}{km} up, about \qty{12}{h}.
    \item \textcolor{hscseriesB}{\bfseries GEO}: \qty{35900}{km} up, 1 day, 6.6 Earth radii from
          the centre.
  \end{itemize}

  \smallskip
  \Cref{ch4:tab:leogeo} links each property to a use. In an HSC answer, always say \emph{which}
  property makes an orbit good for the job.
\end{sidebyside}

\begin{table}
  \caption{Low Earth orbit and geostationary orbit compared.}\label{ch4:tab:leogeo}
  \centering
  \small
  \begin{tabular}{@{}L{28mm}L{56mm}L{56mm}@{}}
    \toprule
    \colhead{Property} & \colhead{Low Earth orbit} & \colhead{Geostationary orbit}\\
    \midrule
    Altitude & \qtyrange{200}{2000}{km} (ISS: \qty{420}{km}) & about \qty{36000}{km}\\
    Speed, period & about \qty{7.7}{km.s^{-1}}, about \qty{90}{min} & \qty{3.1}{km.s^{-1}}, 1 day\\
    In the sky & crosses the sky in minutes & fixed, so ground dishes never move\\
    View & small area, fine detail & about 40\,\% of the Earth at once; poor near the poles\\
    Signal time (up, down) & a few milliseconds & about \qty{0.24}{s}\\
    Problems & orbital decay; many satellites needed for constant cover & weaker signals; costly launch\\
    Uses & ISS, imaging, mapping, polar weather satellites, internet constellations (large groups
      of satellites that work together)
      & TV, communications, weather images of a whole hemisphere (half of the Earth)\\
    \bottomrule
  \end{tabular}
\end{table}

\begin{formulacard}[combine={$\frac{r^3}{T^2} = \frac{GM}{4\pi^2}$}]{Orbit radius and altitude from the period}
  \cardformula{r = \sqrt[3]{\frac{GMT^2}{4\pi^2}} \qquad h = r - r_\mathrm{E}}
  \begin{symbols}
    \sym{$r$}{orbital radius, from the Earth's centre}{\si{m}}
    \sym{$T$}{period of the orbit}{\si{s}}
    \sym{$G$}{universal gravitational constant}{\si{N.m^2.kg^{-2}}}
    \sym{$M$}{mass of the central body (for the Earth, $M_\mathrm{E}$)}{\si{kg}}
    \sym{$h$}{altitude: the height above the surface}{\si{m}}
    \sym{$r_\mathrm{E}$}{radius of the Earth, \num{6.371e6} (data sheet)}{\si{m}}
  \end{symbols}
  \inwords{the period fixes the size of the orbit. Subtract the Earth's radius to get the height.}
  \whentouse{to find how high a satellite must be for a given period (geostationary, GPS, any
    circular orbit).}
  \whennottouse{for an elliptical orbit. There $r$ is the semi-major axis, and the height changes
    around the orbit.}
  \rearranged{$T = 2\pi\sqrt{\dfrac{r^3}{GM}}$ \qquad $r = r_\mathrm{E} + h$}
  \tensecond{$T = \qty{6.00}{h} = \qty{2.16e4}{s}$ gives $r = \qty{1.68e7}{m}$, so
    $h = \qty{1.68e7}{m} - \qty{0.637e7}{m} = \qty{1.04e7}{m}$.}
  \derive{start from Kepler's third law (on the sheet).
    \begin{align*}
      r^3 &= \frac{GMT^2}{4\pi^2} && \text{multiply both sides by $T^2$}\\
      r &= \sqrt[3]{\frac{GMT^2}{4\pi^2}} && \text{take the cube root of both sides}\\
      h &= r - r_\mathrm{E} && \text{subtract the Earth's radius ($r = r_\mathrm{E} + h$)}
    \end{align*}}
  \trap{$T$ must be in seconds (24 h = \qty{86400}{s}). The cube root gives $r$ from the centre.
    Subtract $r_\mathrm{E}$ for the altitude.}
\end{formulacard}

\begin{workedexample}{The geostationary orbit}\label{ch4:we:geo}
Calculate the orbital radius, the altitude and the speed of a geostationary satellite. Use
$T = \qty{24.0}{h}$.
\given{$T = 24.0 \times \qty{3600}{s} = \qty{8.64e4}{s}$, \ $GM_\mathrm{E} = \qty{4.002e14}{m^3.s^{-2}}$, \
  $r_\mathrm{E} = \qty{6.371e6}{m}$}
\find{$r$, the altitude $h$, the speed $v$}
\formula{$r = \sqrt[3]{\dfrac{GMT^2}{4\pi^2}}$ (Kepler's third law), \ $h = r - r_\mathrm{E}$, \
  $v = \dfrac{2\pi r}{T}$}
\substitute{\eqlines{r^3 &= \frac{\qty{4.002e14}{m^3.s^{-2}} \times (\qty{8.64e4}{s})^2}{39.48}
    = \qty{7.567e22}{m^3}\\
  r &= \sqrt[3]{\qty{7.567e22}{m^3}} = \qty{4.230e7}{m} \quad\text{(cube root)}\\
  r_\mathrm{E} &= \qty{6.371e6}{m} = \qty{0.6371e7}{m} \quad\text{(same power of ten as $r$)}\\
  h &= \qty{4.230e7}{m} - \qty{0.6371e7}{m} = \qty{3.593e7}{m}\\
  v &= \frac{2\pi \times \qty{4.230e7}{m}}{\qty{8.64e4}{s}} = \qty{3.08e3}{m.s^{-1}}}}
\answer{\final{$r = \qty{4.23e7}{m}$}, \ \final{$h = \qty{3.59e7}{m}$} (about \qty{36000}{km}), \
  \final{$v = \qty{3.08e3}{m.s^{-1}}$}}
\checkstep{$r$ is 6.6 Earth radii, and the speed is lower than in LEO, as it should be for a
  higher orbit \cmark. With $T$ = 23\,h\,56\,min you get $h = \qty{3.58e7}{m}$; both are accepted
  if you show your $T$.}
\end{workedexample}

\begin{workedexample}{Internet from low Earth orbit}\label{ch4:we:internet}
An internet company uses satellites in circular orbits \qty{550}{km} above the Earth.
(a) Calculate their period in minutes. (b) A signal travels at \qty{3.00e8}{m.s^{-1}} from the
ground to a satellite directly overhead and back. Find the round-trip time for one of these
satellites and for a geostationary one (altitude \qty{3.59e7}{m}). (c) Give one advantage and
one disadvantage of the low orbit.
\given{$r = \qty{6.371e6}{m} + \qty{0.550e6}{m} = \qty{6.921e6}{m}$, \
  $GM_\mathrm{E} = \qty{4.002e14}{m^3.s^{-2}}$, \ $c = \qty{3.00e8}{m.s^{-1}}$ (the speed of light)}
\find{(a) $T$ \quad (b) $t_\mathrm{LEO}$ and $t_\mathrm{GEO}$ \quad (c) one advantage and one
  disadvantage}
\formula{$T = 2\pi\sqrt{\dfrac{r^3}{GM}}$, \quad time $= \dfrac{\text{distance}}{\text{speed}}
  = \dfrac{2h}{c}$ (up and down)}
\substitute{\eqlines{T &= 2\pi\sqrt{\frac{(\qty{6.921e6}{m})^3}{\qty{4.002e14}{m^3.s^{-2}}}}
    = 2\pi \times \qty{910.2}{s} = \qty{5.72e3}{s}\\
  T &= \qty{5.72e3}{s} \div 60 = \qty{95.3}{min}\\
  t_\mathrm{LEO} &= \frac{2 \times \qty{5.50e5}{m}}{\qty{3.00e8}{m.s^{-1}}} = \qty{3.67e-3}{s}\\
  t_\mathrm{GEO} &= \frac{2 \times \qty{3.59e7}{m}}{\qty{3.00e8}{m.s^{-1}}} = \qty{0.239}{s}}}
\answer{(a) \final{$T = \qty{95.3}{min}$} \ (b) \final{$t_\mathrm{LEO} = \qty{3.67e-3}{s}$},
  \final{$t_\mathrm{GEO} = \qty{0.239}{s}$} (65 times longer). (c) Advantage: the delay is tiny,
  so video calls feel instant. Disadvantage: each satellite is overhead for only a few minutes,
  so a large constellation (a group of many satellites) is needed for constant cover.}
\checkstep{the period is a little longer than the ISS's \qty{93}{min}, which fits a slightly
  higher orbit \cmark.}
\end{workedexample}

\begin{mistakes}
\begin{itemize}
  \item \xmark\ Giving $r$ when the question asks for the altitude. \cmark\ $h = r - r_\mathrm{E}$.
  \item \xmark\ ``A geostationary satellite is not moving.'' \cmark\ It moves at
        \qty{3.1}{km.s^{-1}}, keeping pace with the Earth's spin.
  \item \xmark\ ``Any 24-hour orbit is geostationary.'' \cmark\ It must also be circular, above the
        equator and moving west to east.
  \item \xmark\ Listing uses without reasons. \cmark\ Link each use to a property.
\end{itemize}
\end{mistakes}

\begin{examtip}
``Compare low Earth and geostationary orbits and relate them to their uses'' (4--5 marks). For
each orbit, give a property with a number (altitude or period), then the use it allows, joined
by ``so'' or ``because''. Use a table or ``whereas'' to compare. ``Compare'' means similarities
\emph{and} differences, so also give one similarity. For example, in both orbits gravity provides
the centripetal force, and both obey $\dfrac{r^3}{T^2} = \dfrac{GM}{4\pi^2}$ with the Earth's mass.
In calculations, write $T$ in seconds (\qty{86400}{s}), take a cube root, then subtract
$r_\mathrm{E}$.
\end{examtip}

\begin{quickcheck}
  \item State the three conditions for a geostationary orbit.
  \item Some navigation satellites have a period of \qty{14.1}{h}. Calculate their altitude.
  \item A \qty{24}{h} orbit is tilted at \ang{40} to the equator. Is it geostationary? Explain.
  \item Explain why a satellite that photographs single buildings must be in a low orbit.
\end{quickcheck}

%% ==================================================================
\section{Key experiments}\label{ch4:sec:experiments}
\covers{M5.IQ3.01, M5.IQ3.02, M5.IQ3.04}

The dot points say \emph{investigate}. A school cannot change the mass of a planet, so the first
two investigations use \term{secondary data}: measurements made by others, such as space
agencies. The third uses an orbit simulation, where you can change one thing at a time.
(\Cref{ch5:sec:experiment} uses a simulation to study energy in elliptical orbits.)

\begin{experiment}{What affects the field strength \texorpdfstring{$g$}{g}?}
\aim{To use secondary data to show that a body's surface field strength is proportional to
  $\dfrac{M}{r^2}$ ($M$: mass, $r$: radius of the body), as $g = \dfrac{GM}{r^2}$ predicts.}
\apparatus{a reliable data source (such as a space agency's planetary fact sheet) and a
  spreadsheet (\cref{ch4:fig:exp1flow}).
  \begin{figure}
    \centering
    \begin{tikzpicture}[node distance=8mm]
      \node[concept=33mm] (a) {data source: $M$, $r$ and surface $g$ for 5 bodies};
      \node[concept=30mm, right=of a] (b) {spreadsheet: calculate $M/r^2$};
      \node[concept=33mm, right=of b] (c) {graph of $g$ against $M/r^2$; find the gradient};
      \draw[link] (a) -- (b);
      \draw[link] (b) -- (c);
    \end{tikzpicture}
    \caption{The steps of this secondary-data investigation.}\label{ch4:fig:exp1flow}
  \end{figure}}
\method{\begin{enumerate}[beginpenalty=10000]
  \item Record $M$, $r$ and the surface $g$ for the Moon, Mercury, Mars, Venus and the Earth.
  \item Calculate $\dfrac{M}{r^2}$ for each body.
  \item Plot $g$ (vertical) against $\dfrac{M}{r^2}$ (horizontal). Draw the line of best fit and
        find its gradient.
\end{enumerate}}
\results{{\centering\small % not a center environment: it would let the page break after the heading
  \smallskip
  \begin{tabular}{@{}lcccc@{}}
    \toprule
    Body & $M$ (\si{kg}) & $r$ (\si{m}) & $M/r^2$ ($\times10^{10}$ \si{kg.m^{-2}}) & $g$ (\si{m.s^{-2}})\\
    \midrule
    Moon    & \num{7.35e22} & \num{1.737e6} & 2.44  & 1.62\\
    Mercury & \num{3.30e23} & \num{2.440e6} & 5.54  & 3.70\\
    Mars    & \num{6.42e23} & \num{3.390e6} & 5.59  & 3.71\\
    Venus   & \num{4.87e24} & \num{6.052e6} & 13.30 & 8.87\\
    Earth   & \num{5.97e24} & \num{6.371e6} & 14.71 & 9.80\\
    \bottomrule
  \end{tabular}\par}
  \medskip
  In the $M/r^2$ column, 2.44 means $2.44 \times 10^{10}$ \si{kg.m^{-2}}. These values are more
  precise than the data sheet, so the Earth's mass is \qty{5.97e24}{kg} here.
  \begin{figure}
  \centering
  \begin{tikzpicture}
    \begin{axis}[hsc, width=0.62\linewidth, height=0.34\linewidth,
        xmin=0, xmax=16, ymin=0, ymax=11,
        xlabel={$M/r^2$ ($\times10^{10}$ \si{kg.m^{-2}})}, ylabel={$g$ (\si{m.s^{-2}})},
        xtick={0,4,8,12,16}, ytick={0,2,4,6,8,10}]
      \addplot[hscseriesB, dashed, line width=1pt, domain=0:16, samples=2, mark=none] {0.6665*x};
      \addplot[hscseriesA, only marks, mark=*, mark size=2pt]
        coordinates {(2.436,1.62) (5.543,3.70) (5.586,3.71) (13.30,8.87) (14.71,9.80)};
      \node[anchor=west, font=\small] at (axis cs:7.0,1.6) {gradient $= \qty{6.67e-11}{N.m^2.kg^{-2}}$};
    \end{axis}
  \end{tikzpicture}
  \caption{Surface field strength against $M/r^2$ for five bodies.}\label{ch4:fig:exp1graph}
  \end{figure}}
\conclusion{The points lie on a straight line through the origin, so $g \propto \dfrac{M}{r^2}$.
  The gradient, \qty{6.67e-11}{N.m^2.kg^{-2}}, equals $G$. This is consistent with
  $g = \dfrac{GM}{r^2}$.}
\significance{Used in reverse, $M = \dfrac{gr^2}{G}$ gives a planet's mass from $g$ and $r$, once
  $G$ is known. In 1798 Henry Cavendish measured the tiny pull between lead spheres. His result is
  the same as measuring $G$, so it gave the Earth's mass for the first time (``weighing the
  Earth'').}
\validity{\begin{itemize}[beginpenalty=10000]% keep the heading with the first item
  \item \textbf{Validity:} use one trusted source, so all bodies are measured the same way (for
        example, the average radius of each body).
  \item \textbf{Limitation (validity):} space agencies find a planet's mass from $GM$ (measured
        from the orbits of moons or spacecraft) divided by $G$. Many tables also calculate surface
        $g$ from $\dfrac{GM}{r^2}$. So the data are not fully independent of the law. The graph
        checks that the data agree with the law, but it is not a real test of it. The Earth's $g$
        can be measured directly (for example with falling objects or a pendulum). For most other
        bodies, $g$ is calculated.
  \item \textbf{Reliability:} check the data against a second, independent source. Use five or
        more bodies and the gradient of the line, not one point.
  \item \textbf{Accuracy:} compare the gradient with $G$ on the data sheet. The line of best fit
        gives \num{6.665e-11}, which differs from \num{6.67e-11} by less than 0.1\,\%.
\end{itemize}}
\end{experiment}

\begin{experiment}{Kepler's third law and the moons of Jupiter}
\aim{To show that $\dfrac{r^3}{T^2}$ is the same for all the moons of one planet, and to use it to
  find the mass of Jupiter.}
\apparatus{a telescope (or published observatory data) and a spreadsheet. Galileo first saw these
  four moons in 1610 (\cref{ch4:fig:jupiter}).
  \begin{figure}
    \centering
    \begin{tikzpicture}
      \foreach \y/\lab/\pos in {0/Night 1/{-2.6/C,-1.2/E,0.9/I,2.0/G},
                                -1.0/Night 2/{-1.9/G,-0.7/I,1.4/E,3.2/C}} {
        \node[diagramlabel, anchor=east] at (-3.7,\y) {\lab};
        \fill[hscbrand!12, draw=hscink, line width=0.6pt] (0,\y) circle (0.26);
        \foreach \x/\m in \pos {
          \fill[hscink] (\x,\y) circle (1.6pt);
          \node[diagramlabel, above, inner sep=3pt, font=\footnotesize] at (\x,\y) {\m};
        }
      }
      \node[diagramlabel] at (0,0.55) {Jupiter};
      \draw[guide] (0,-1.26) -- (0,-1.6);
      \draw[guide] (3.2,-1.05) -- (3.2,-1.6);
      \draw[hscmuted, line width=0.6pt, {Stealth[length=2mm]}-{Stealth[length=2mm]}]
        (0,-1.5) -- node[below, diagramlabel] {$d$} (3.2,-1.5);
    \end{tikzpicture}
    \caption{Jupiter's four large moons (I = Io, E = Europa, G = Ganymede, C = Callisto) in a
      small telescope on two nights (not to scale). They move along a line as they orbit. $d$ is a
      moon's distance from Jupiter on one night. Its largest value, over many nights, is the
      orbital radius $r$.}\label{ch4:fig:jupiter}
  \end{figure}}
\method{\begin{enumerate}[beginpenalty=10000]
  \item For each moon, record its largest distance from Jupiter over many nights (as an angle, or
        in Jupiter diameters). Use Jupiter's known diameter, \qty{1.43e8}{m}, to change this
        distance to metres. This distance is the orbital radius $r$. (Published data already give
        $r$.) Also record the period $T$: the time to return to the same position.
  \item Change $T$ to seconds. Calculate $r^3$ and $T^2$.
  \item Plot $r^3$ (vertical) against $T^2$ (horizontal) and find the gradient. Then
        $M = \dfrac{4\pi^2 \times \text{gradient}}{G}$.
\end{enumerate}}
\results{{\centering\small % not a center environment (see the first experiment)
  \smallskip
  \begin{tabular}{@{}lcccc@{}}
    \toprule
    Moon & $r$ (\si{m}) & $T$ (\si{s}) & $r^3$ ($\times10^{27}$ \si{m^3}) & $T^2$ ($\times10^{12}$ \si{s^2})\\
    \midrule
    Io       & \num{4.22e8} & \num{1.53e5} & 0.0752 & 0.0234\\
    Europa   & \num{6.71e8} & \num{3.07e5} & 0.302  & 0.0942\\
    Ganymede & \num{1.07e9} & \num{6.18e5} & 1.23   & 0.382\\
    Callisto & \num{1.88e9} & \num{1.44e6} & 6.64   & 2.07\\
    \bottomrule
  \end{tabular}\par}
  \medskip
  In the $r^3$ column, 0.0752 means $0.0752 \times 10^{27}$ \si{m^3}. In the $T^2$ column, 0.0234
  means $0.0234 \times 10^{12}$ \si{s^2}.
  \begin{figure}
  \centering
  \begin{tikzpicture}
    \begin{axis}[hsc, width=0.62\linewidth, height=0.34\linewidth,
        xmin=0, xmax=2.3, ymin=0, ymax=7.5,
        xlabel={$T^2$ ($\times10^{12}$ \si{s^2})}, ylabel={$r^3$ ($\times10^{27}$ \si{m^3})},
        xtick={0,0.5,1,1.5,2}, ytick={0,2,4,6},
        xticklabel style={/pgf/number format/.cd, fixed, precision=1}]
      \addplot[hscseriesB, dashed, line width=1pt, domain=0:2.3, samples=2, mark=none] {3.2045*x};
      \addplot[hscseriesA, only marks, mark=*, mark size=2pt]
        coordinates {(0.02341,0.07515) (0.09425,0.3021) (0.3819,1.225) (2.074,6.645)};
      \node[anchor=west, font=\small] at (axis cs:0.9,1.0) {gradient $= \qty{3.20e15}{m^3.s^{-2}}$};
    \end{axis}
  \end{tikzpicture}
  \caption{Kepler's third law for the moons of Jupiter.}\label{ch4:fig:exp2graph}
  \end{figure}}
\conclusion{The graph is a straight line through the origin, so $r^3 \propto T^2$ (Kepler's
  third law). The gradient is $\dfrac{GM}{4\pi^2} = \qty{3.205e15}{m^3.s^{-2}}$, so
  \[ M = \frac{4\pi^2 \times \text{gradient}}{G} = \frac{39.48 \times \num{3.205e15}}{\num{6.67e-11}}
     = \qty{1.90e27}{kg} \]
  This is the mass of Jupiter (accepted value \qty{1.90e27}{kg}).}
\significance{Jupiter's moons were the first objects found orbiting another planet. This showed that not
  everything orbits the Earth, which supported the Sun-centred model. The same method gives the
  masses of planets, stars and black holes.}
\validity{\begin{itemize}[beginpenalty=10000]
  \item \textbf{Validity:} all four moons orbit the same planet, so they share one value of
        $\dfrac{GM}{4\pi^2}$. Their orbits are nearly circular and their masses are tiny compared
        with Jupiter's, as the formula assumes.
  \item \textbf{Reliability:} time many orbits, not one. Use the gradient of the line of best fit,
        not a single moon.
  \item \textbf{Accuracy:} the mass agrees with the accepted value to better than 0.1\,\%. Change
        all units to metres and seconds first.
\end{itemize}}
\end{experiment}

% layout QA: short title on the continuation part, so the title bar stays on one line
\begin{experiment}[title after break={\hscboxtitle{\faFlask}{Key experiment}{Orbits in a simulation}\hspace{0.5em}{\footnotesize\mdseries\itshape(continued)}}]{Orbits in a simulation: mass, radius, speed and period}
\aim{To find how a satellite's orbital period and speed depend on its own mass and on its
  orbital radius.}
\apparatus{\begin{sidebyside}[0.40]
    \centering
    \begin{tikzpicture}
      \fill[hscbrand!10, draw=hscink, line width=0.8pt] (0,0) circle (0.5);
      \node[diagramlabel, font=\footnotesize] at (0,0) {Earth};
      \draw[hscmuted, dashed, line width=0.6pt] (0,0) circle (1.4);
      \draw[guide] (0,0) -- (1.4,0);
      \node[diagramlabel, below, font=\footnotesize] at (0.95,0) {$r$};
      \draw[velocity] (1.4,0) -- (1.4,1.0) node[right, text=hscvel, font=\footnotesize] {$v$};
      \fill[object] (1.4,0) circle (3.5pt);
      \node[diagramlabel, right, font=\footnotesize, inner sep=5pt] at (1.4,-0.1) {$m$};
    \end{tikzpicture}
    \captionof{figure}{Simulation set-up: launch at \ang{90} to the radius.}\label{ch4:fig:simsetup}
  \nextside
    A computer orbit simulation (for example PhET ``Gravity and Orbits'', or a Python model) with
    a timer (\cref{ch4:fig:simsetup}). Set the planet's mass to the Earth's, \qty{6.0e24}{kg}.
    Launch the satellite at \ang{90} to the radius. Adjust the launch speed until the path is a
    circle.
  \end{sidebyside}}
\method{\begin{enumerate}[beginpenalty=10000]
  \item Set $r = \qty{7.0e6}{m}$. Time 5 orbits and divide by 5 to get the period $T$.
  \item Keep $r$ the same. Change only the satellite's mass (\qty{100}{kg}, \qty{1000}{kg},
        \qty{5000}{kg}). Find $T$ each time.
  \item Keep the mass the same. Repeat for $r$ = \num{1.0e7}, \num{2.0e7} and \qty{4.0e7}{m}.
        For each radius calculate $v = \dfrac{2\pi r}{T}$ and then $v^2 r$.
\end{enumerate}}
\results{The values below come from a simple Python model of the same set-up. Changing the mass
  (\qty{100}{kg}, \qty{1000}{kg} or \qty{5000}{kg}) at $r = \qty{7.0e6}{m}$ gave the same period
  every time, $T = \qty{5817}{s}$. For different radii:
  \begin{center}\small
  \begin{tabular}{@{}cccc@{}}
    \toprule
    $r$ (\si{m}) & $T$ (\si{s}) & $v = \dfrac{2\pi r}{T}$ (\si{m.s^{-1}}) & $v^2 r$ (\si{m^3.s^{-2}})\\
    \midrule
    \num{7.0e6} & 5817  & 7561 & \num{4.00e14}\\
    \num{1.0e7} & 9932  & 6326 & \num{4.00e14}\\
    \num{2.0e7} & 28092 & 4473 & \num{4.00e14}\\
    \num{4.0e7} & 79457 & 3163 & \num{4.00e14}\\
    \bottomrule
  \end{tabular}
  \end{center}}
\conclusion{The period does not depend on the satellite's mass. A bigger orbit has a longer
  period and a lower speed. $v^2 r$ is the same for every radius, so $v^2 \propto \dfrac{1}{r}$
  and $v \propto \dfrac{1}{\sqrt r}$. This is what $v = \sqrt{\dfrac{GM}{r}}$ predicts, because it
  gives $v^2 r = GM$.}
\significance{All objects at one height move together, whatever their mass. That is why an
  astronaut floats beside the space station. Engineers choose a satellite's height, not its mass,
  to set its period.}
\validity{\begin{itemize}[beginpenalty=10000]
  \item \textbf{Validity:} change only one variable at a time (first the mass, then the radius).
        The simulation is built on Newton's law, so it shows what the law predicts. It does not
        test the law against real orbits.
  \item \textbf{Reliability:} time 5 orbits, not one. Repeat each run with a smaller time step;
        the results should not change.
  \item \textbf{Accuracy:} compare $v^2 r$ with
        $GM_\mathrm{E} = \num{6.67e-11} \times \num{6.0e24} = \qty{4.002e14}{m^3.s^{-2}}$
        (percentage difference). Here they agree to better than 0.1\,\%.
\end{itemize}}
\end{experiment}

%% ==================================================================
% Chapter summary (added by the whole-book consistency edit: "How to use this guide" says
% that every chapter ends with a short summary). Every statement is taught above.
\begin{summarybox}[unbreakable]{\Cref{ch4:chap} in five lines}
\begin{itemize}
  \item $F = \frac{GMm}{r^2}$ and $g = \frac{GM}{r^2}$, with $r = r_\mathrm{E} + h$ measured from
        the centre of the planet. Double $r$, and $F$ and $g$ become 4 times smaller.
  \item Gravity provides the centripetal force, so $v = \sqrt{\frac{GM}{r}}$. The satellite's mass
        cancels, and a higher orbit is slower.
  \item Kepler's third law: $\frac{r^3}{T^2} = \frac{GM}{4\pi^2}$ is the same for everything that
        orbits one central body. Put $T$ in seconds, and use a cube root to find $r$.
  \item Kepler's first two laws: each planet's orbit is an ellipse with the Sun at one focus,
        and equal areas are swept in equal times, so a planet is fastest at perihelion.
  \item Low Earth orbit (\qtyrange{200}{2000}{km} up, about \qty{90}{min}): fine detail and a
        short signal delay. Geostationary orbit (about \qty{36000}{km} up, one day, above the equator): it
        stays above one point.
\end{itemize}
\end{summarybox}
